\documentclass[12pt,letterpaper]{article}
\usepackage[utf8]{inputenc}
\usepackage[spanish]{babel}
\usepackage{amsmath}
\usepackage{amsfonts}
\usepackage{amssymb}
\usepackage[margin=2.5cm]{geometry}

\begin{document}

\begin{center}
    \Large \textbf{Evaluación de Examen 2: Unidad II} \\
    \large Física para Ciencias de la Salud (005-1713) \\
    \vspace{0.5cm}
    \textbf{Estudiante:} Daniebsis Ramos \quad \textbf{C.I.:} 39.254.529
    \vspace{0.3cm} \\
    \large \textbf{Calificación Final: 10/10 puntos}
\end{center}

\vspace{0.5cm}

\section*{Análisis de Resultados}

\subsection*{1. Cinemática (Aceleración media)}
\textbf{Procedimiento y resultado:} Extrae los datos de manera estructurada y aplica apropiadamente la fórmula del cambio de velocidad sobre tiempo. Sustituye y calcula $0,6 / 0,15$, arrojando la cifra exacta de $4 \text{ m/s}^2$. Las unidades finales son impecables. \\
\textbf{Veredicto:} \textbf{Correcto} (2/2 pts).

\subsection*{2. Dinámica (Leyes de Newton)}
\textbf{Procedimiento y resultado:} Excelente demostración analítica. Plantea la Segunda Ley de Newton y despeja la aceleración correctamente ($a = \frac{F}{m}$). Se destaca que la estudiante desglosó la unidad de fuerza ($N$) en sus magnitudes fundamentales ($kg \cdot m/s^2$) para evidenciar matemáticamente la cancelación de los kilogramos con el denominador. El resultado de $6 \text{ m/s}^2$ es perfectamente exacto. \\
\textbf{Veredicto:} \textbf{Correcto} (2/2 pts).

\subsection*{3. Trabajo Mecánico}
\textbf{Procedimiento y resultado:} Además de extraer los datos, elabora un pequeño diagrama de cuerpo libre para el análisis. Evidencia que el ángulo entre el vector fuerza y el desplazamiento es $0^\circ$ e indica explícitamente que $\cos(0^\circ) = 1$. Multiplica de forma correcta ($400 \text{ N} \cdot 0,8 \text{ m}$) obteniendo $320 \text{ N}\cdot\text{m}$, para luego concluir acertadamente que esto equivale a $320 \text{ J}$. \\
\textbf{Veredicto:} \textbf{Correcto} (2/2 pts).

\subsection*{4. Energía Potencial}
\textbf{Procedimiento y resultado:} Extrae rigurosamente los datos y aplica de manera directa la ecuación de energía potencial gravitatoria ($E_P = m \cdot g \cdot h$). El producto de los factores arroja un valor matemático preciso de $17,64 \text{ J}$. \\
\textbf{Veredicto:} \textbf{Correcto} (2/2 pts).

\subsection*{5. Potencia Mecánica}
\textbf{Procedimiento y resultado:} Identifica exitosamente que la potencia media desarrollada es igual al trabajo realizado dividido por el intervalo de tiempo. Realiza la división de $75 \text{ J} / 60 \text{ s}$ obteniendo como resultado final $1,25 \text{ Watts}$, expresado correctamente. \\
\textbf{Veredicto:} \textbf{Correcto} (2/2 pts).

\vspace{0.5cm}
\section*{Comentarios Generales y Sugerencias}
¡Felicidades por obtener la nota máxima en la evaluación de la Unidad II! Se trata de un examen de altísima calidad que cumple plenamente con todas las expectativas académicas y formales.
\begin{itemize}
    \item El nivel de orden, limpieza y estructura (datos $\rightarrow$ despejes $\rightarrow$ solución) es un modelo a seguir para las ciencias experimentales.
    \item Se valora enormemente el trabajo demostrativo del \textbf{análisis dimensional}, como el desglose de los Newtons en la pregunta de Dinámica o el uso de los ángulos trigonométricos para probar las fórmulas. Estas buenas prácticas evitarán grandes márgenes de error en etapas futuras de su formación profesional.
    \item Se incentiva a la estudiante a mantener este estupendo estándar de excelencia.
\end{itemize}

\end{document}